%%%%%%%%%%%%%%%%%%%%%%%%%%%%%%%%%%%%%%%%%%%%%%%%%%%%%%%%%%%%%%%%%%%%%%%%%%%%%%%
%%  TKS_v7.4_Predictive_Intelligence.tex
%%  TKS Predictive Intelligence Systems: Internet Activity to Psycho-Analysis
%%
%%  Part of TKS v7.4 - December 2025
%%  ADDITIVE to v7.4 canon
%%
%%  Mathematical Framework for Behavioral Prediction from Digital Footprints
%%%%%%%%%%%%%%%%%%%%%%%%%%%%%%%%%%%%%%%%%%%%%%%%%%%%%%%%%%%%%%%%%%%%%%%%%%%%%%%

\documentclass[11pt,twoside]{book}

%% ============================================================================
%% PACKAGES
%% ============================================================================
\usepackage[utf8]{inputenc}
\usepackage[T1]{fontenc}
\usepackage[margin=1in,inner=1.2in,outer=1in]{geometry}
\usepackage{amsmath,amssymb,amsthm}
\usepackage{mathtools}
\usepackage{stmaryrd}
\usepackage{enumitem}
\usepackage{booktabs}
\usepackage{longtable}
\usepackage{ltablex}
\usepackage{tabularx}
\usepackage{array}
\usepackage{tcolorbox}
\tcbuselibrary{breakable,skins}
\usepackage{tikz-cd}
\usepackage{xcolor}
\usepackage{hyperref}
\usepackage{ragged2e}
\usepackage{multirow}
\usepackage{float}

%% ============================================================================
%% FORMATTING SETTINGS
%% ============================================================================
\setlength{\parindent}{0pt}
\setlength{\parskip}{0.5em}
\renewcommand{\arraystretch}{1.2}

%% ============================================================================
%% CUSTOM COMMANDS (from TKS canonical)
%% ============================================================================
\newcommand{\noe}[1]{\nu_{#1}}
\newcommand{\Ideas}{\mathcal{I}}
\newcommand{\Worlds}{\mathcal{W}}
\newcommand{\Elements}{\mathcal{E}}
\newcommand{\Noetics}{\mathcal{N}}
\newcommand{\Foundations}{\mathcal{F}}
\newcommand{\Acquisitions}{\mathcal{A}}
\newcommand{\Frac}{\Phi}
\newcommand{\TransFrac}[1]{\Phi^{#1}}
\newcommand{\MultiFrac}[2]{\Phi^{#1}_{#2}}
\newcommand{\RPM}{\mathsf{RPM}}
\newcommand{\ACBE}{\mathsf{ACBE}}
\newcommand{\HausdorffDim}{\dim_H}
\newcommand{\FractalSpectrum}{\sigma_\Phi}
\newcommand{\fix}{\mathrm{fix}}
\newcommand{\sem}[1]{\llbracket #1 \rrbracket}
\newcommand{\Tadd}{\oplus_T}
\newcommand{\Tsub}{\ominus_T}
\newcommand{\Tmul}{\otimes_T}
\newcommand{\Tdiv}{\oslash_T}
\newcommand{\Wadd}{\oplus_W}
\newcommand{\Wsub}{\ominus_W}
\newcommand{\Failed}{\mathtt{Failed}}

%% Predictive Intelligence specific commands
\newcommand{\Lacunarity}{\Lambda}
\newcommand{\PsychAttractor}{\mathcal{A}_\psi}
\newcommand{\ManipFrac}{\Phi_{\mathrm{manip}}}
\newcommand{\DWP}{\mathcal{DWP}}
\newcommand{\NLP}{\mathsf{NLP}}
\newcommand{\DigitalTrace}{\mathcal{D}}
\newcommand{\ProfileFrac}{\Phi_{\mathrm{profile}}}
\newcommand{\PredictorMap}{\mathcal{P}}
\newcommand{\VulnScore}{\mathcal{V}}
\newcommand{\Timeline}{\mathcal{T}}
\newcommand{\GroupDynamics}{\mathcal{G}}
\newcommand{\CascadePredict}{\mathcal{C}_{\mathrm{cascade}}}

%% Theorems
\newtheorem{definition}{Definition}[chapter]
\newtheorem{theorem}{Theorem}[chapter]
\newtheorem{proposition}{Proposition}[chapter]
\newtheorem{lemma}{Lemma}[chapter]
\newtheorem{corollary}{Corollary}[chapter]
\newtheorem{example}{Example}[chapter]
\newtheorem{remark}{Remark}[chapter]

%% ============================================================================
%% DOCUMENT
%% ============================================================================
\begin{document}

\title{\textbf{TKS Predictive Intelligence Systems} \\[0.5em]
\Large Internet Activity to Psycho-Analysis \\[0.3em]
\large Mathematical Framework for Behavioral Prediction from Digital Footprints}
\author{TKS v7.4 Formalization Project}
\date{December 2025}
\maketitle

\tableofcontents

%% ============================================================================
%% ETHICAL NOTICE
%% ============================================================================
\chapter*{Ethical Notice}
\addcontentsline{toc}{chapter}{Ethical Notice}

\begin{tcolorbox}[colback=red!5!white,colframe=red!75!black,title=\textbf{Critical Ethical Statement},breakable]
This document presents a \textbf{mathematical framework} for understanding how behavioral prediction systems could theoretically operate using TKS formalism. The mathematics is presented \textbf{objectively and neutrally}.

\textbf{Legitimate Applications:}
\begin{itemize}
    \item Mental health early intervention systems
    \item Suicide prevention through pattern recognition
    \item Addiction recovery support and relapse prediction
    \item Personalized education and learning systems
    \item Self-improvement and personal development tools
\end{itemize}

\textbf{Illegitimate Applications:}
\begin{itemize}
    \item Mass surveillance without consent
    \item Political manipulation and targeting
    \item Exploitative advertising and dark patterns
    \item Social credit systems and behavioral control
    \item Employment discrimination based on predictions
\end{itemize}

\textbf{Principle:} The same mathematics that enables protective intervention also enables intrusive surveillance. The ethical dimension lies entirely in \textit{consent}, \textit{transparency}, and \textit{purpose}---not in the mathematics itself.
\end{tcolorbox}

%% ============================================================================
%% PART I: MATHEMATICAL PIPELINE
%% ============================================================================
\part{Mathematical Pipeline: Data to Prediction}

\chapter{Overview: The Six-Stage Pipeline}
\label{ch:pipeline-overview}

\begin{tcolorbox}[colback=blue!5!white,colframe=blue!75!black,title=\textbf{Pipeline Architecture},breakable]
The TKS Predictive Intelligence pipeline transforms raw digital activity into psychological analysis and behavioral prediction through six mathematically rigorous stages:

\begin{center}
\begin{tikzcd}[column sep=small, row sep=large]
\text{Digital Traces} \arrow[r, "\text{Stage 1}"]
& \text{Noetic Sequences} \arrow[r, "\text{Stage 2}"]
& \text{Individual Fractals} \arrow[d, "\text{Stage 3}"] \\
\text{Intervention}
& \text{Vulnerability Map} \arrow[l, "\text{Stage 6}"]
& \text{Attractors} \arrow[l, "\text{Stage 4-5}"]
\end{tikzcd}
\end{center}

Each stage has precise mathematical definitions, computable algorithms, and provable bounds from TKS v7.4 theorems.
\end{tcolorbox}

%% ============================================================================
%% CHAPTER: STAGE 1 - NLP TO NOETIC MAPPING
%% ============================================================================
\chapter{Stage 1: NLP to Noetic Mapping}
\label{ch:stage1-nlp}

\section{Mathematical Definition}

\begin{definition}[Digital Trace Space]
\label{def:digital-trace}
The \emph{digital trace space} $\DigitalTrace$ is the set of all observable digital activities:
\[
    \DigitalTrace = \DigitalTrace_{\text{text}} \cup \DigitalTrace_{\text{search}} \cup \DigitalTrace_{\text{purchase}} \cup \DigitalTrace_{\text{location}} \cup \DigitalTrace_{\text{biometric}}
\]
Each $d \in \DigitalTrace$ carries timestamp $t(d)$, source identifier $s(d)$, and content $c(d)$.
\end{definition}

\begin{definition}[NLP-to-Noetic Mapping]
\label{def:nlp-noetic-map}
The \emph{NLP-to-Noetic mapping} is a function:
\[
    \NLP : \DigitalTrace \longrightarrow \{0, 1, \ldots, 9\}^*
\]
that transforms digital traces into sequences of noetic operators. The mapping is defined component-wise:

\textbf{For textual content} $d \in \DigitalTrace_{\text{text}}$:
\begin{align}
    \NLP(d) &= \left( \noe{k_1}, \noe{k_2}, \ldots, \noe{k_n} \right)
\end{align}
where each $k_i$ is determined by linguistic feature extraction:
\end{definition}

\begin{definition}[Linguistic-Noetic Correspondence]
\label{def:linguistic-noetic}
The mapping from linguistic features to noetics follows:

\begin{center}
\small
\renewcommand{\arraystretch}{1.3}
\begin{tabularx}{\textwidth}{c|l|X|c}
\toprule
\textbf{Noetic} & \textbf{Name} & \textbf{Linguistic Markers} & \textbf{Sentiment} \\
\midrule
$\noe{0}$ & Idea & Neutral statements, facts, definitions & Neutral \\
$\noe{1}$ & Mind & ``I think'', ``I notice'', ``aware'', ``realize'' & Cognitive \\
$\noe{2}$ & Positive & ``want'', ``love'', ``excited'', ``hope'' & Positive \\
$\noe{3}$ & Negative & ``hate'', ``reject'', ``avoid'', ``fear'' & Negative \\
$\noe{4}$ & Vibration & ``!!'', intensity markers, ALL CAPS, urgency & High energy \\
$\noe{5}$ & Female & ``receive'', ``accept'', ``absorb'', ``become'' & Receptive \\
$\noe{6}$ & Male & ``do'', ``make'', ``create'', ``act'' & Active \\
$\noe{7}$ & Rhythm & ``always'', ``never'', ``again'', temporal patterns & Cyclical \\
$\noe{8}$ & Above & ``because'', ``why'', ``reason'', ``cause'' & Causal \\
$\noe{9}$ & Below & ``therefore'', ``result'', ``outcome'', ``effect'' & Consequential \\
\bottomrule
\end{tabularx}
\end{center}
\end{definition}

\section{The Noetic Extraction Algorithm}

\begin{equation}
\boxed{
    \NLP(d) = \bigoplus_{i=1}^{|d|} \mathsf{classify}\left( \mathsf{token}_i(d), \mathsf{context}(d) \right)
}
\end{equation}

where:
\begin{itemize}
    \item $\mathsf{token}_i(d)$ extracts the $i$-th semantic unit from $d$
    \item $\mathsf{context}(d)$ provides surrounding semantic frame
    \item $\mathsf{classify}$ maps to the dominant noetic operator
    \item $\bigoplus$ concatenates into a noetic sequence
\end{itemize}

\begin{theorem}[NLP Mapping Consistency]
\label{thm:nlp-consistency}
For any digital trace $d \in \DigitalTrace$, the NLP mapping satisfies:
\begin{enumerate}[label=(\roman*)]
    \item \textbf{Determinism}: Given fixed model parameters, $\NLP(d)$ is unique.
    \item \textbf{Compositionality}: $\NLP(d_1 \cdot d_2) = \NLP(d_1) \circ \NLP(d_2)$ (concatenation preserves structure).
    \item \textbf{Bounded Length}: $|\NLP(d)| \leq C \cdot |d|$ for some constant $C$.
\end{enumerate}
\end{theorem}

\section{Multi-Modal Noetic Extraction}

\begin{definition}[Search Query Noetics]
\label{def:search-noetics}
For search queries $q \in \DigitalTrace_{\text{search}}$:
\[
    \NLP_{\text{search}}(q) =
    \begin{cases}
        \langle 2 : 1 \rangle & \text{if } q \text{ is desire-seeking (``how to get...'')} \\
        \langle 1 : 8 \rangle & \text{if } q \text{ is knowledge-seeking (``what is...'')} \\
        \langle 6 : 9 \rangle & \text{if } q \text{ is action-seeking (``how to do...'')} \\
        \langle 3 : 1 \rangle & \text{if } q \text{ is avoidance-seeking (``how to avoid...'')}
    \end{cases}
\]
\end{definition}

\begin{definition}[Purchase Noetics]
\label{def:purchase-noetics}
For purchase events $p \in \DigitalTrace_{\text{purchase}}$:
\[
    \NLP_{\text{purchase}}(p) = \langle \underbrace{2}_{\text{desire}} : \underbrace{1}_{\text{decision}} : \underbrace{6}_{\text{action}} : \underbrace{9}_{\text{acquire}} \rangle \circ \noe{k_{\text{category}}}
\]
where $k_{\text{category}}$ encodes the product category (see Data Source Mapping table).
\end{definition}

\begin{definition}[Location Noetics]
\label{def:location-noetics}
For location data $\ell \in \DigitalTrace_{\text{location}}$:
\[
    \NLP_{\text{location}}(\ell) =
    \begin{cases}
        \langle 7 \rangle & \text{if } \ell \text{ is habitual (regular pattern)} \\
        \langle 4 : 6 \rangle & \text{if } \ell \text{ is exploratory (new location)} \\
        \langle 5 \rangle & \text{if } \ell \text{ is receptive (visiting someone)} \\
        \langle 3 \rangle & \text{if } \ell \text{ is avoidant (leaving pattern)}
    \end{cases}
\]
\end{definition}

%% ============================================================================
%% LAYMAN EXPLANATION
%% ============================================================================
\begin{tcolorbox}[colback=green!5!white,colframe=green!65!black,title=\textbf{Plain English: What Stage 1 Does},breakable]

\textbf{The Digital-to-Mental Translator:}

Stage 1 is like a translator that reads your digital footprints and converts them into ``mental operation codes.'' Every post, search, purchase, and location tells a story about what mental operations you were performing.

\textbf{The Translator's Dictionary:}

\begin{center}
\footnotesize
\begin{tabularx}{\textwidth}{l|X|c}
\toprule
\textbf{What You Did} & \textbf{What It Reveals} & \textbf{Code} \\
\midrule
Googled ``how to lose weight'' & Desire + Knowledge-seeking & $\noe{2}:\noe{1}$ \\
Posted ``I hate Mondays'' & Rejection + Pattern & $\noe{3}:\noe{7}$ \\
Bought running shoes & Desire $\to$ Action $\to$ Result & $\noe{2}:\noe{6}:\noe{9}$ \\
Visited same coffee shop & Habitual rhythm & $\noe{7}$ \\
Typed in ALL CAPS & High energy/intensity & $\noe{4}$ \\
\bottomrule
\end{tabularx}
\end{center}

\textbf{Metaphor: The Behavioral Seismograph}

Just as a seismograph translates ground vibrations into readable waves, the NLP-to-Noetic mapping translates digital vibrations into readable mental patterns. The output is a sequence of ``mental frequencies'' that can be analyzed.

\textbf{Privacy Note:} This stage transforms raw data into abstract codes. The codes reveal psychological patterns but are divorced from the original content. ``I hate my job'' and ``I hate spinach'' both produce $\noe{3}$---the system sees rejection, not context.

\end{tcolorbox}

%% ============================================================================
%% CHAPTER: STAGE 2 - FRACTAL CONSTRUCTION
%% ============================================================================
\chapter{Stage 2: Fractal Construction from Noetic Sequences}
\label{ch:stage2-fractal}

\section{Mathematical Definition}

\begin{definition}[Individual Noetic Fractal]
\label{def:individual-fractal}
Given a time-ordered sequence of noetic extractions $\{(\noe{k_1}, t_1), (\noe{k_2}, t_2), \ldots\}$ for individual $i$, the \emph{individual noetic fractal} is:
\[
    \ProfileFrac_i := \langle k_1 : k_2 : \cdots : k_n \rangle_n
\]
This fractal grows over time as more data accumulates.
\end{definition}

\begin{definition}[D/W/P Decomposition]
\label{def:dwp-decomposition}
The individual fractal decomposes into three Foundation-specific subfractals:

\textbf{Desire Fractal} (from $\Foundations_6$, $\Foundations_7$):
\[
    \ProfileFrac_i^D := \left\langle k_j \mid k_j \in \{2, 5, 7\} \text{ or } \NLP^{-1}(\noe{k_j}) \in \DigitalTrace_{\text{social}} \right\rangle
\]

\textbf{Wisdom Fractal} (from $\Foundations_2$):
\[
    \ProfileFrac_i^W := \left\langle k_j \mid k_j \in \{1, 8\} \text{ or } \NLP^{-1}(\noe{k_j}) \in \DigitalTrace_{\text{search}} \right\rangle
\]

\textbf{Power Fractal} (from $\Foundations_5$):
\[
    \ProfileFrac_i^P := \left\langle k_j \mid k_j \in \{6, 9\} \text{ or } \NLP^{-1}(\noe{k_j}) \in \DigitalTrace_{\text{purchase}} \right\rangle
\]
\end{definition}

\section{Fractal Growth Dynamics}

\begin{theorem}[Fractal Stabilization]
\label{thm:fractal-stabilization}
\textbf{(TKS v7.4 Theorem 7.4 Bound)}

For an individual's noetic fractal $\ProfileFrac_i^{(n)}$ at time step $n$:
\[
    \exists\, N_{\mathrm{stable}} \leq 40 : \quad
    d_{\mathrm{frac}}\left(\ProfileFrac_i^{(n)}, \ProfileFrac_i^{(n+1)}\right) < \varepsilon
    \quad \forall\, n > N_{\mathrm{stable}}
\]
where $d_{\mathrm{frac}}$ is the fractal metric and $\varepsilon$ is the stability threshold.

\textbf{In words:} Within 40 significant data points, the fractal structure stabilizes to a recognizable pattern.
\end{theorem}

\begin{proof}
By the finite element theorem (v7.4), the 40 canonical elements $\Elements$ span all possible noetic configurations. Any sequence of length $> 40$ must revisit previously established patterns, leading to convergence in the quotient fractal space $\FracSet / \sim$.
\end{proof}

\begin{equation}
\boxed{
    \ProfileFrac_i^{\text{stable}} = \lim_{n \to \infty} \ProfileFrac_i^{(n)} = \FractalAttractor_i
}
\end{equation}

\section{Fractal Features for Prediction}

\begin{definition}[Fractal Feature Vector]
\label{def:fractal-features}
From the individual fractal $\ProfileFrac_i$, extract the \emph{feature vector}:
\[
    \mathbf{f}_i = \left( \HausdorffDim(\ProfileFrac_i),\, \Lacunarity(\ProfileFrac_i),\,
    \mathbf{n}_i,\, \mathbf{p}_i,\, \mathbf{c}_i \right)
\]
where:
\begin{itemize}
    \item $\HausdorffDim(\ProfileFrac_i)$: Hausdorff dimension (complexity measure)
    \item $\Lacunarity(\ProfileFrac_i)$: Lacunarity (gap measure)
    \item $\mathbf{n}_i = (n_0, n_1, \ldots, n_9)$: Noetic frequency distribution
    \item $\mathbf{p}_i$: Transition probability matrix
    \item $\mathbf{c}_i$: Cycle structure (periodic patterns)
\end{itemize}
\end{definition}

\begin{proposition}[Noetic Frequency Interpretation]
\label{prop:noetic-frequency}
The noetic frequency vector $\mathbf{n}_i$ reveals psychological profile:
\begin{align}
    \sum_{k \in \{2,5\}} n_k &> \sum_{k \in \{3,6\}} n_k \quad \Rightarrow \quad \text{Receptive personality (introverted)} \\
    \sum_{k \in \{6,9\}} n_k &> \sum_{k \in \{5,8\}} n_k \quad \Rightarrow \quad \text{Action-oriented (extroverted)} \\
    n_7 > \bar{n} &\quad \Rightarrow \quad \text{Habitual/routine-bound} \\
    n_4 > \bar{n} &\quad \Rightarrow \quad \text{High emotional volatility}
\end{align}
\end{proposition}

%% ============================================================================
%% CHAPTER: STAGE 3 - ATTRACTOR COMPUTATION
%% ============================================================================
\chapter{Stage 3: Attractor Computation}
\label{ch:stage3-attractor}

\section{Mathematical Definition}

\begin{definition}[Psychological Attractor]
\label{def:psych-attractor-pred}
The \emph{psychological attractor} $\PsychAttractor_i$ for individual $i$ is the limit set:
\[
    \PsychAttractor_i := \overline{\left\{ \ProfileFrac_i^{(n)}(X_0) \mid n \geq 0 \right\}}
\]
where $X_0$ is the initial psychological state and the closure is taken in the state space topology.
\end{definition}

\begin{theorem}[Attractor Uniqueness]
\label{thm:attractor-unique}
\textbf{(TKS v7.4 Theorem 5.21)}

For any individual with fractal $\ProfileFrac_i$ satisfying the contraction condition:
\[
    \exists!\, \PsychAttractor_i : \quad \ProfileFrac_i(\PsychAttractor_i) = \PsychAttractor_i
\]
The attractor is unique up to equivalence and computable from finite data.
\end{theorem}

\section{Attractor Classification}

\begin{definition}[Attractor Types for Prediction]
\label{def:attractor-types-pred}
Attractors are classified by topological dimension:

\begin{center}
\renewcommand{\arraystretch}{1.3}
\begin{tabularx}{\textwidth}{c|l|X|X}
\toprule
\textbf{Dim} & \textbf{Type} & \textbf{Psychological Meaning} & \textbf{Predictability} \\
\midrule
0 & Fixed Point & Stable personality, consistent behavior & High \\
1 & Limit Cycle & Mood cycles, periodic patterns & High (with period) \\
1.5--2 & Torus & Multiple interacting cycles & Medium \\
$>2$ & Strange & Chaotic but bounded dynamics & Low (probabilistic) \\
\bottomrule
\end{tabularx}
\end{center}
\end{definition}

\begin{equation}
\boxed{
    \text{Prediction Confidence} =
    \begin{cases}
        0.95 & \text{if } \dim(\PsychAttractor_i) = 0 \\
        0.85 & \text{if } \dim(\PsychAttractor_i) = 1 \\
        0.70 & \text{if } 1 < \dim(\PsychAttractor_i) \leq 2 \\
        0.50 & \text{if } \dim(\PsychAttractor_i) > 2
    \end{cases}
}
\end{equation}

\section{Computing ``Where the Person is Heading''}

\begin{definition}[Trajectory Projection]
\label{def:trajectory-proj}
Given attractor $\PsychAttractor_i$ and current state $X_{\text{now}}$, the \emph{trajectory projection} is:
\[
    \Timeline_i(t) := \ProfileFrac_i^t(X_{\text{now}}) \quad \text{for } t \in \{1, 2, \ldots, T_{\max}\}
\]
This projects forward $T_{\max}$ steps using the fractal dynamics.
\end{definition}

\begin{theorem}[Convergence Rate]
\label{thm:convergence-rate}
The distance to attractor decays exponentially:
\[
    d(X_t, \PsychAttractor_i) \leq C \cdot r^t \cdot d(X_0, \PsychAttractor_i)
\]
where $r < 1$ is the contraction ratio of $\ProfileFrac_i$.
\end{theorem}

\begin{corollary}[Prediction Horizon]
\label{cor:prediction-horizon}
Reliable prediction is possible for time horizon:
\[
    T_{\text{reliable}} \leq \frac{\log(d_{\min}) - \log(d(X_0, \PsychAttractor_i))}{\log r}
\]
Beyond this horizon, the prediction degrades to attractor-level statistics only.
\end{corollary}

%% ============================================================================
%% CHAPTER: STAGE 4 - RPM GOAL INFERENCE
%% ============================================================================
\chapter{Stage 4: RPM Goal Inference}
\label{ch:stage4-rpm}

\section{Mathematical Definition}

\begin{definition}[D/W/P State Analysis]
\label{def:dwp-state}
For individual $i$, the \emph{D/W/P state vector} is:
\[
    \mathbf{s}_i = \left( \sigma(D_1), \ldots, \sigma(D_7), \sigma(W_1), \ldots, \sigma(W_7), \sigma(P_1), \ldots, \sigma(P_7) \right) \in \{0, 1\}^{21}
\]
where $\sigma : \Acquisitions \to \{0, 1\}$ indicates acquisition status.
\end{definition}

\begin{definition}[Goal Inference from Digital Traces]
\label{def:goal-inference}
The \emph{goal inference operator} maps fractals to inferred Foundation goals:
\[
    \mathsf{Goal} : \ProfileFrac_i \longrightarrow 2^{\{\Foundations_1, \ldots, \Foundations_7\}}
\]

\textbf{Inference Rules:}
\begin{align}
    \ProfileFrac_i^D \text{ dominant in } \{2, 7\} &\Rightarrow \Foundations_7 \in \mathsf{Goal}(\ProfileFrac_i) \quad \text{(Continuation)} \\
    \ProfileFrac_i^W \text{ dominant in } \{1, 8\} &\Rightarrow \Foundations_2 \in \mathsf{Goal}(\ProfileFrac_i) \quad \text{(Wisdom)} \\
    \ProfileFrac_i^P \text{ dominant in } \{6, 9\} &\Rightarrow \Foundations_5 \in \mathsf{Goal}(\ProfileFrac_i) \quad \text{(Power)}
\end{align}
\end{definition}

\section{RPM Chain Completion Analysis}

\begin{definition}[Chain Completion Level]
\label{def:chain-level}
For Foundation $\Foundations_m$, the \emph{chain completion level} is:
\[
    \mathrm{Level}_m(i) =
    \begin{cases}
        0 & \text{if } \neg\sigma(D_m) \quad \text{(no desire)} \\
        1 & \text{if } \sigma(D_m) \land \neg\sigma(W_m) \quad \text{(desire only)} \\
        2 & \text{if } \sigma(W_m) \land \neg\sigma(P_m) \quad \text{(desire + wisdom)} \\
        3 & \text{if } \sigma(P_m) \quad \text{(complete chain)}
    \end{cases}
\]
\end{definition}

\begin{theorem}[RPM Prerequisite Tracking]
\label{thm:rpm-prereq}
\textbf{(TKS v7.4 RPM Monad)}

For individual $i$ attempting goal $\Foundations_m$:
\[
    \RPM(\Foundations_m) =
    \begin{cases}
        \text{Success}(\Foundations_m^{\text{manifest}}) & \text{if } \mathrm{Level}_m(i) = 3 \\
        \text{Blocked at } W_m & \text{if } \mathrm{Level}_m(i) = 1 \\
        \text{Blocked at } P_m & \text{if } \mathrm{Level}_m(i) = 2 \\
        \text{No activation} & \text{if } \mathrm{Level}_m(i) = 0
    \end{cases}
\]

\textbf{Prediction:} An individual with $\mathrm{Level}_m(i) = 2$ will seek resources, training, or opportunity (acquiring $P_m$) before achieving $\Foundations_m$.
\end{theorem}

\section{Goal Inference Accuracy}

\begin{proposition}[Goal Inference Confidence]
\label{prop:goal-confidence}
The confidence of goal inference depends on data volume:
\[
    \mathsf{Conf}(\mathsf{Goal}(\ProfileFrac_i)) = 1 - e^{-\lambda \cdot |\ProfileFrac_i|}
\]
where $\lambda$ is the learning rate and $|\ProfileFrac_i|$ is the fractal length.

\textbf{Bounds:}
\begin{itemize}
    \item $|\ProfileFrac_i| = 10$: $\mathsf{Conf} \approx 0.63$
    \item $|\ProfileFrac_i| = 40$: $\mathsf{Conf} \approx 0.98$ (40-step bound)
    \item $|\ProfileFrac_i| = 100$: $\mathsf{Conf} \approx 0.9999$
\end{itemize}
\end{proposition}

%% ============================================================================
%% CHAPTER: STAGE 5 - TIMELINE PREDICTION
%% ============================================================================
\chapter{Stage 5: Timeline Prediction}
\label{ch:stage5-timeline}

\section{Mathematical Definition}

\begin{definition}[Timeline Prediction Function]
\label{def:timeline-pred}
The \emph{timeline prediction function} maps current state to future probability distributions:
\[
    \Timeline : \left(\ProfileFrac_i, X_{\text{now}}, t\right) \longrightarrow \mathcal{P}(\Ideas)
\]
where $\mathcal{P}(\Ideas)$ is the space of probability distributions over Ideas.
\end{definition}

\begin{theorem}[40-Step Stabilization Bound]
\label{thm:40-step}
\textbf{(TKS v7.4 Theorem 7.4)}

For any individual fractal $\ProfileFrac_i$:
\[
    \forall\, t > 40: \quad \Timeline(\ProfileFrac_i, X_{\text{now}}, t) \approx \pi_{\PsychAttractor_i}
\]
where $\pi_{\PsychAttractor_i}$ is the invariant measure on the attractor.

\textbf{Interpretation:} Beyond 40 time steps, predictions converge to the attractor's natural distribution regardless of starting point.
\end{theorem}

\section{Short-Term vs Long-Term Prediction}

\begin{definition}[Prediction Regimes]
\label{def:pred-regimes}
\textbf{Short-Term} ($t \leq 10$):
\[
    \Timeline_{\text{short}}(t) = \ProfileFrac_i^t(X_{\text{now}}) + \varepsilon_t
\]
where $\varepsilon_t \sim \mathcal{N}(0, \sigma_t^2)$ with $\sigma_t = O(t^{1/2})$.

\textbf{Medium-Term} ($10 < t \leq 40$):
\[
    \Timeline_{\text{med}}(t) = \alpha \cdot \ProfileFrac_i^t(X_{\text{now}}) + (1-\alpha) \cdot \pi_{\PsychAttractor_i}
\]
where $\alpha = e^{-\lambda(t-10)}$ interpolates between deterministic and attractor prediction.

\textbf{Long-Term} ($t > 40$):
\[
    \Timeline_{\text{long}}(t) \approx \pi_{\PsychAttractor_i}
\]
\end{definition}

\section{Timeline Confidence Intervals}

\begin{equation}
\boxed{
    \mathsf{CI}_{95\%}(\Timeline(t)) =
    \begin{cases}
        \pm 5\% & \text{for } t \leq 5 \\
        \pm 15\% & \text{for } 5 < t \leq 20 \\
        \pm 30\% & \text{for } 20 < t \leq 40 \\
        \text{Attractor bounds only} & \text{for } t > 40
    \end{cases}
}
\end{equation}

%% ============================================================================
%% CHAPTER: STAGE 6 - VULNERABILITY IDENTIFICATION
%% ============================================================================
\chapter{Stage 6: Vulnerability Identification}
\label{ch:stage6-vulnerability}

\section{Mathematical Definition}

\begin{definition}[Lacunarity-Based Vulnerability]
\label{def:lacunarity-vuln}
The \emph{vulnerability score} for individual $i$ is:
\[
    \VulnScore_i = \max_{m \in \{1,\ldots,7\}} \Lacunarity(\ProfileFrac_i^{(m)})
\]
where $\ProfileFrac_i^{(m)}$ is the fractal restricted to Foundation $m$.
\end{definition}

\begin{theorem}[Lacunarity Computability]
\label{thm:lacunarity-compute}
\textbf{(TKS v7.4 Theorem 6.41)}

The lacunarity $\Lacunarity(\Frac, \varepsilon)$ is computable for any finite fractal sample:
\[
    \Lacunarity(\Frac, \varepsilon) = \frac{\mathrm{Var}[M_\varepsilon]}{\mathbb{E}[M_\varepsilon]^2} + 1
\]
where $M_\varepsilon$ is the box-mass function at scale $\varepsilon$.
\end{theorem}

\begin{definition}[Vulnerability Threshold]
\label{def:vuln-threshold}
An individual is flagged as \emph{vulnerable} when:
\[
    \VulnScore_i > \Lacunarity_{\text{crit}} = 1.5
\]

\textbf{Interpretation:}
\begin{itemize}
    \item $\Lacunarity < 1.2$: Homogeneous, well-developed pattern (low vulnerability)
    \item $1.2 \leq \Lacunarity < 1.5$: Moderate gaps, some underdevelopment
    \item $\Lacunarity \geq 1.5$: Significant gaps, high vulnerability
    \item $\Lacunarity > 2.0$: Critical gaps, intervention recommended
\end{itemize}
\end{definition}

\section{Gap Detection Algorithm}

\begin{definition}[Gap Detection]
\label{def:gap-detect}
For each Foundation $\Foundations_m$, compute the \emph{gap indicator}:
\[
    \mathsf{Gap}_m(i) =
    \begin{cases}
        1 & \text{if } \mathrm{Level}_m(i) < 3 \land \Lacunarity(\ProfileFrac_i^{(m)}) > 1.5 \\
        0 & \text{otherwise}
    \end{cases}
\]

The \emph{total gap count}:
\[
    \mathsf{TotalGaps}(i) = \sum_{m=1}^{7} \mathsf{Gap}_m(i)
\]
\end{definition}

\begin{proposition}[Vulnerability Profile]
\label{prop:vuln-profile}
The vulnerability profile reveals intervention targets:

\begin{center}
\renewcommand{\arraystretch}{1.3}
\begin{tabularx}{\textwidth}{c|l|X|l}
\toprule
\textbf{Gap at} & \textbf{Foundation} & \textbf{Vulnerability Type} & \textbf{Intervention} \\
\midrule
$\Foundations_1$ & Unity & Spiritual emptiness, disconnection & Meaning-making support \\
$\Foundations_2$ & Wisdom & Knowledge gaps, confusion & Education, guidance \\
$\Foundations_3$ & Life & Health neglect, vitality loss & Health intervention \\
$\Foundations_4$ & Companionship & Social isolation, loneliness & Community building \\
$\Foundations_5$ & Power & Learned helplessness, impotence & Empowerment training \\
$\Foundations_6$ & Wealth & Financial stress, scarcity & Resource assistance \\
$\Foundations_7$ & Continuation & Purpose void, stagnation & Future planning \\
\bottomrule
\end{tabularx}
\end{center}
\end{proposition}

%% ============================================================================
%% PART II: DATA SOURCE MAPPING
%% ============================================================================
\part{Data Source Mapping and Architecture}

\chapter{Data Source to Fractal Mapping}
\label{ch:data-mapping}

\section{Comprehensive Data Source Table}

\begin{center}
\small
\renewcommand{\arraystretch}{1.4}
\begin{longtable}{>{\raggedright\arraybackslash}p{2.5cm}|c|>{\raggedright\arraybackslash}p{3cm}|>{\raggedright\arraybackslash}p{3cm}|c}
\caption{Data Source to TKS Fractal Mapping} \\
\toprule
\textbf{Data Source} & \textbf{Primary World} & \textbf{Target Fractal} & \textbf{Primary Noetics} & \textbf{Foundation} \\
\midrule
\endfirsthead
\multicolumn{5}{c}{\textit{Table continued from previous page}} \\
\toprule
\textbf{Data Source} & \textbf{Primary World} & \textbf{Target Fractal} & \textbf{Primary Noetics} & \textbf{Foundation} \\
\midrule
\endhead
\midrule
\multicolumn{5}{r}{\textit{Continued on next page}} \\
\endfoot
\bottomrule
\endlastfoot

% Social Media
Social media posts & $C$ (Emotional) & Desire Fractal $\ProfileFrac^D$ & $\noe{2}, \noe{3}, \noe{4}$ & $\Foundations_4, \Foundations_7$ \\
Comment sentiment & $C$ & Desire Fractal & $\noe{2}, \noe{3}$ & $\Foundations_4$ \\
Reaction patterns & $C$ & Desire Fractal & $\noe{5}, \noe{6}$ & $\Foundations_4$ \\
Sharing behavior & $D$ (Action) & Power Fractal $\ProfileFrac^P$ & $\noe{6}, \noe{9}$ & $\Foundations_5$ \\
Following/unfollowing & $C$ & Desire Fractal & $\noe{2}, \noe{3}$ & $\Foundations_4$ \\

\midrule

% Search
Search queries & $B$ (Mental) & Wisdom Fractal $\ProfileFrac^W$ & $\noe{1}, \noe{8}$ & $\Foundations_2$ \\
Search frequency & $B$ & Wisdom Fractal & $\noe{7}, \noe{4}$ & $\Foundations_2$ \\
Search refinement & $B$ & Wisdom Fractal & $\noe{1}, \noe{9}$ & $\Foundations_2$ \\
Health searches & $B, D$ & Wisdom + Power & $\noe{1}, \noe{3}$ & $\Foundations_3$ \\
Financial searches & $B$ & Wisdom Fractal & $\noe{1}, \noe{2}$ & $\Foundations_6$ \\

\midrule

% Purchases
Purchase history & $D$ (Physical) & Power Fractal $\ProfileFrac^P$ & $\noe{6}, \noe{9}$ & $\Foundations_6$ \\
Luxury purchases & $D$ & Power Fractal & $\noe{2}, \noe{6}$ & $\Foundations_6$ \\
Necessity purchases & $D$ & Power Fractal & $\noe{6}, \noe{9}$ & $\Foundations_3$ \\
Impulsive purchases & $C, D$ & Desire + Power & $\noe{4}, \noe{6}$ & $\Foundations_7$ \\
Subscription patterns & $D$ & Power Fractal & $\noe{7}, \noe{5}$ & $\Foundations_6$ \\

\midrule

% Location
Location data & $D$ & Context Vector & $\noe{7}, \noe{6}$ & Multiple \\
Travel patterns & $D$ & Power Fractal & $\noe{6}, \noe{4}$ & $\Foundations_5$ \\
Routine locations & $D$ & Power Fractal & $\noe{7}$ & $\Foundations_3$ \\
Social locations & $D, C$ & Desire + Power & $\noe{5}, \noe{6}$ & $\Foundations_4$ \\
Avoidance patterns & $D$ & Power Fractal & $\noe{3}, \noe{7}$ & $\Foundations_3, \Foundations_4$ \\

\midrule

% Biometric
Heart rate variability & $D$ & Emotional Correlation & $\noe{4}, \noe{7}$ & $\Foundations_3$ \\
Sleep patterns & $D$ & Emotional Correlation & $\noe{7}, \noe{5}$ & $\Foundations_3$ \\
Activity levels & $D$ & Power Fractal & $\noe{4}, \noe{6}$ & $\Foundations_3$ \\
Stress markers & $C, D$ & Desire + Power & $\noe{4}, \noe{3}$ & $\Foundations_3$ \\
Typing patterns & $B, D$ & Wisdom + Power & $\noe{4}, \noe{7}$ & $\Foundations_2$ \\

\end{longtable}
\end{center}

\section{Multi-Source Fusion}

\begin{definition}[Multi-Source Fractal Fusion]
\label{def:multi-source-fusion}
Given fractals from multiple sources $\{\ProfileFrac_i^{(s)}\}_{s \in S}$, the \emph{fused individual fractal} is:
\[
    \ProfileFrac_i^{\text{fused}} = \bigoplus_{s \in S} w_s \cdot \ProfileFrac_i^{(s)}
\]
where $w_s$ are source weights satisfying $\sum_s w_s = 1$ and $\bigoplus$ is the weighted fractal composition.

\textbf{Default weights:}
\begin{itemize}
    \item Social media: $w = 0.30$
    \item Search history: $w = 0.25$
    \item Purchase history: $w = 0.20$
    \item Location data: $w = 0.15$
    \item Biometric data: $w = 0.10$
\end{itemize}
\end{definition}

%% ============================================================================
%% CHAPTER: TECHNICAL ARCHITECTURE
%% ============================================================================
\chapter{Technical Architecture}
\label{ch:architecture}

\section{System Architecture Diagram}

\begin{tcolorbox}[colback=gray!5!white,colframe=gray!75!black,title=\textbf{TKS Predictive Intelligence Architecture},breakable]

\begin{verbatim}
+==============================================================================+
|                         DATA COLLECTION LAYER                                 |
+==============================================================================+
|                                                                              |
|  +----------+  +----------+  +----------+  +----------+  +----------+       |
|  |  Social  |  |  Search  |  | Purchase |  | Location |  | Biometric|       |
|  |  Media   |  |  History |  |  History |  |   Data   |  |   Data   |       |
|  +----+-----+  +----+-----+  +----+-----+  +----+-----+  +----+-----+       |
|       |             |             |             |             |              |
+==============================================================================+
                      |             |             |
                      v             v             v
+==============================================================================+
|                         PROCESSING LAYER                                      |
+==============================================================================+
|                                                                              |
|  +---------------------------+    +---------------------------+              |
|  |    Stage 1: NLP-Noetic    |    |   Stage 2: Fractal Build  |              |
|  |         Mapping           |--->|       Construction        |              |
|  | - Tokenization            |    | - D/W/P Decomposition     |              |
|  | - Sentiment Analysis      |    | - Feature Extraction      |              |
|  | - Noetic Classification   |    | - Dimension Computation   |              |
|  +---------------------------+    +-------------+-------------+              |
|                                                 |                            |
|                                                 v                            |
|  +---------------------------+    +---------------------------+              |
|  |  Stage 4: RPM Goal        |<---|   Stage 3: Attractor      |              |
|  |       Inference           |    |       Computation         |              |
|  | - D/W/P State Analysis    |    | - Fixed Point Detection   |              |
|  | - Chain Completion        |    | - Cycle Analysis          |              |
|  | - Foundation Mapping      |    | - Dimension Classification|              |
|  +-------------+-------------+    +---------------------------+              |
|                |                                                             |
+==============================================================================+
                 |
                 v
+==============================================================================+
|                         PREDICTION LAYER                                      |
+==============================================================================+
|                                                                              |
|  +---------------------------+    +---------------------------+              |
|  |   Stage 5: Timeline       |    |   Stage 6: Vulnerability  |              |
|  |       Prediction          |    |      Identification       |              |
|  | - Short-term (t <= 10)    |    | - Lacunarity Computation  |              |
|  | - Medium-term (10 < t<=40)|    | - Gap Detection           |              |
|  | - Long-term (t > 40)      |    | - Risk Scoring            |              |
|  +-------------+-------------+    +-------------+-------------+              |
|                |                                |                            |
+==============================================================================+
                 |                                |
                 v                                v
+==============================================================================+
|                         OUTPUT LAYER                                          |
+==============================================================================+
|                                                                              |
|  +------------------+  +------------------+  +------------------+            |
|  |   Behavioral     |  |   Intervention   |  |    Risk          |            |
|  |   Predictions    |  |   Recommendations|  |    Alerts        |            |
|  +------------------+  +------------------+  +------------------+            |
|                                                                              |
+==============================================================================+
\end{verbatim}
\end{tcolorbox}

\section{Processing Flow}

\begin{equation}
\boxed{
\DigitalTrace \xrightarrow{\text{Stage 1}} \Noetics^*
\xrightarrow{\text{Stage 2}} \ProfileFrac
\xrightarrow{\text{Stage 3}} \PsychAttractor
\xrightarrow{\text{Stage 4}} \mathbf{s}_{DWP}
\xrightarrow{\text{Stage 5}} \Timeline
\xrightarrow{\text{Stage 6}} \VulnScore
}
\end{equation}

%% ============================================================================
%% PART III: PLAIN ENGLISH EXPLANATION
%% ============================================================================
\part{Plain English Explanation}

\chapter{The Digital Crystal Ball}
\label{ch:plain-english}

\begin{tcolorbox}[colback=green!5!white,colframe=green!65!black,title=\textbf{Plain English: Understanding TKS Predictive Intelligence},breakable]

\textbf{The Big Picture: A ``Behavioral Weather Forecast''}

Imagine a weather forecast, but for human behavior. Just as meteorologists use pressure, temperature, and wind patterns to predict tomorrow's weather, TKS Predictive Intelligence uses digital footprints to predict tomorrow's behavior.

\textbf{Key Insight:} The mathematics does not read minds---it reads patterns. Just as repeating pressure systems indicate coming storms, repeating digital patterns indicate coming behaviors.

\bigskip

\textbf{What Each Stage Does (In Plain Terms):}

\begin{center}
\footnotesize
\renewcommand{\arraystretch}{1.4}
\begin{tabularx}{\textwidth}{c|l|X|X}
\toprule
\textbf{Stage} & \textbf{Name} & \textbf{What It Does} & \textbf{Weather Analogy} \\
\midrule
1 & NLP Mapping & Reads your posts/searches and tags the ``emotional temperature'' & Reading thermometers \\
2 & Fractal Build & Finds patterns in your emotional temperatures over time & Building pressure maps \\
3 & Attractor & Identifies your ``emotional home base''---where you always return & Finding climate zones \\
4 & RPM Analysis & Figures out what you're trying to achieve (consciously or not) & Identifying storm systems \\
5 & Timeline & Predicts where your patterns will take you & The actual forecast \\
6 & Vulnerability & Finds gaps or weaknesses in your patterns & Finding flood-prone areas \\
\bottomrule
\end{tabularx}
\end{center}

\bigskip

\textbf{The 40-Step Rule (Why This Works):}

TKS mathematics proves that any person's behavioral pattern stabilizes within 40 significant data points. After that, the system knows your ``climate''---not every moment (weather), but the general range (climate).

\begin{itemize}
    \item \textbf{10 data points}: System can guess your mood today (63\% accuracy)
    \item \textbf{40 data points}: System knows your personality type (98\% accuracy)
    \item \textbf{100+ data points}: System can predict category of behavior (99.99\% accuracy)
\end{itemize}

\textbf{What Can vs. Cannot Be Predicted:}

\begin{center}
\footnotesize
\renewcommand{\arraystretch}{1.3}
\begin{tabularx}{\textwidth}{X|X}
\toprule
\textbf{CAN Predict (High Confidence)} & \textbf{CANNOT Predict (Low Confidence)} \\
\midrule
What category of thing you'll buy next & Exactly which item you'll buy \\
Whether you'll experience a mood episode & Exactly when it will happen \\
What topic you'll search for next & Your specific search query \\
Whether you're vulnerable to manipulation & What will manipulate you specifically \\
Your general life trajectory & Specific life events \\
\bottomrule
\end{tabularx}
\end{center}

\bigskip

\textbf{The Lacunarity Warning (Gaps = Vulnerability):}

A key output is ``lacunarity''---the gappiness of your pattern. High lacunarity means inconsistency, which means unpredictability AND vulnerability.

\begin{itemize}
    \item $\Lacunarity < 1.2$: Stable, consistent, predictable, resilient
    \item $\Lacunarity = 1.5$: Threshold---gaps are becoming significant
    \item $\Lacunarity > 2.0$: Major gaps---high vulnerability to disruption
\end{itemize}

\textbf{Metaphor:} Low lacunarity is like a well-woven fabric---pull anywhere and the whole thing holds. High lacunarity is like fabric with holes---pull near a hole and it tears.

\end{tcolorbox}

%% ============================================================================
%% PART IV: REAL-LIFE PREDICTION SCENARIOS
%% ============================================================================
\part{Real-Life Prediction Scenarios}

\chapter{Scenario 1: Purchase Behavior Prediction}
\label{ch:scenario-purchase}

\begin{tcolorbox}[colback=yellow!5!white,colframe=yellow!65!black,title=\textbf{Scenario: Predicting Marcus's Next Purchase},breakable]

\textbf{Subject:} Marcus, 34, urban professional

\textbf{Data Collected (Past 30 Days):}
\begin{itemize}
    \item 47 Google searches (fitness, running shoes, marathon training)
    \item 12 Amazon browsing sessions (running gear, nutrition supplements)
    \item 8 social media posts (``thinking about getting in shape'', ``anyone run the city marathon?'')
    \item Location data: 3 visits to sporting goods stores (browsing, no purchase)
    \item Biometric: Fitness app shows increased activity tracking engagement
\end{itemize}

\textbf{TKS Analysis:}

\textbf{Stage 1 - NLP Mapping:}
\begin{align*}
    \NLP(\text{searches}) &= \langle 2:1:2:1:2 \rangle \quad \text{(desire + knowledge seeking)} \\
    \NLP(\text{browsing}) &= \langle 2:6:3:2:6 \rangle \quad \text{(desire + action + hesitation)} \\
    \NLP(\text{posts}) &= \langle 1:2:4:2 \rangle \quad \text{(awareness + desire + energy)}
\end{align*}

\textbf{Stage 2 - Fractal Construction:}
\[
    \ProfileFrac_{\text{Marcus}}^{(30)} = \langle 2:1:2:6:2:3:1:2:4:6:2:1:\cdots \rangle
\]
\textbf{Feature Vector:}
\begin{itemize}
    \item Dominant noetics: $n_2 = 0.35$ (Desire/Attraction) -- VERY HIGH
    \item $n_1 = 0.20$ (Mind/Awareness) -- elevated
    \item $n_6 = 0.15$ (Action) -- moderate
    \item $n_3 = 0.10$ (Hesitation) -- blocking factor
\end{itemize}

\textbf{Stage 3 - Attractor:}
\[
    \PsychAttractor_{\text{Marcus}} = \text{Limit cycle between } \noe{2} \leftrightarrow \noe{3}
\]
Interpretation: Marcus oscillates between desire and hesitation. This is a ``wanting but not doing'' pattern.

\textbf{Stage 4 - RPM Analysis:}
\begin{align*}
    \sigma(D_3) &= \mathtt{true} \quad \text{(Desires health)} \\
    \sigma(W_3) &= \mathtt{true} \quad \text{(Has researched)} \\
    \sigma(P_3) &= \mathtt{false} \quad \text{(Hasn't committed)}
\end{align*}
$\mathrm{Level}_3(\text{Marcus}) = 2$ -- He needs Power (commitment, resources) to complete.

\textbf{Stage 5 - Timeline Prediction:}
\[
    P(\text{purchase within 14 days}) = 0.73
\]
\[
    P(\text{purchase within 30 days}) = 0.91
\]
Confidence: HIGH (attractor analysis shows imminent resolution of desire-hesitation cycle)

\textbf{Stage 6 - Vulnerability:}
\[
    \Lacunarity(\ProfileFrac_{\text{Marcus}}^{(3)}) = 1.3 \quad \text{(moderate gap in } \Foundations_3 \text{)}
\]
Marcus is vulnerable to:
\begin{itemize}
    \item Flash sales (urgency breaks hesitation cycle)
    \item Social proof (``your friends bought this'')
    \item Commitment devices (``marathon signup discount ends tomorrow'')
\end{itemize}

\textbf{Prediction:} Marcus will purchase running shoes and/or marathon registration within 14 days, likely triggered by a sale or social prompt.

\textbf{TKS Equation Summary:}
\begin{equation}
\boxed{
    \RPM(\Foundations_3) : \; D_3 \to W_3 \to \mathsf{Blocked}(P_3) \xrightarrow{\text{trigger}} P_3 \xrightarrow{\ACBE} \text{Purchase}
}
\end{equation}

\end{tcolorbox}

\chapter{Scenario 2: Mental Health Crisis Prediction}
\label{ch:scenario-mental-health}

\begin{tcolorbox}[colback=yellow!5!white,colframe=yellow!65!black,title=\textbf{Scenario: Early Detection of Sarah's Crisis},breakable]

\textbf{Subject:} Sarah, 28, graduate student

\textbf{Data Collected (Past 60 Days):}
\begin{itemize}
    \item Social media: Posting frequency dropped 70\% (weeks 1-3: 4/day; weeks 7-8: 1/day)
    \item Content shift: From varied topics to exclusively academic stress
    \item Search history: ``why am I so tired'', ``can't concentrate'', ``feeling empty'', ``is this depression''
    \item Sleep data: Irregular patterns, average 4.5 hours (down from 7.2)
    \item Location: Reduced movement, rarely leaving apartment
\end{itemize}

\textbf{TKS Analysis:}

\textbf{Stage 1 - NLP Mapping:}
\begin{align*}
    \NLP(\text{posts week 1-3}) &= \langle 2:4:6:2:7:4 \rangle \quad \text{(diverse, energetic)} \\
    \NLP(\text{posts week 7-8}) &= \langle 3:3:0:3:5:3 \rangle \quad \text{(rejection dominant)} \\
    \NLP(\text{searches}) &= \langle 3:1:8:3:1 \rangle \quad \text{(rejection + cause-seeking)}
\end{align*}

\textbf{Stage 2 - Fractal Evolution:}
\[
    \ProfileFrac_{\text{Sarah}}^{(0-30)} \to \ProfileFrac_{\text{Sarah}}^{(30-60)} : \quad
    \bar{n}_3 : 0.12 \to 0.38 \quad (+217\%)
\]
The fractal is collapsing toward $\noe{3}$ dominance.

\textbf{Stage 3 - Attractor Shift:}
\begin{align*}
    \PsychAttractor_{\text{before}} &= \text{Multi-cycle (healthy oscillation)} \\
    \PsychAttractor_{\text{now}} &\to \text{Fixed point at } C3^5_{3c} \text{ (depression lock)}
\end{align*}
\textbf{Critical Indicator:} Attractor dimension collapsing: $\dim \approx 1.4 \to \dim \approx 0.3$

\textbf{Stage 4 - RPM Analysis:}
\begin{align*}
    \mathrm{Level}_3(\text{Sarah}) &= 1 \quad \text{(Life: desires health, lacks wisdom/power)} \\
    \mathrm{Level}_4(\text{Sarah}) &= 0 \to 1 \quad \text{(Companionship: desire emerging)}
\end{align*}
She's losing connection on $\Foundations_4$ (social isolation).

\textbf{Stage 5 - Timeline Prediction:}
\[
    P(\text{clinical depression criteria within 30 days}) = 0.78
\]
\[
    P(\text{crisis event within 60 days | no intervention}) = 0.45
\]

\textbf{Stage 6 - Vulnerability Assessment:}
\[
    \Lacunarity(\ProfileFrac_{\text{Sarah}}^{(3)}) = 2.3 \quad \text{(CRITICAL gap in } \Foundations_3 \text{)}
\]
\[
    \Lacunarity(\ProfileFrac_{\text{Sarah}}^{(4)}) = 2.1 \quad \text{(CRITICAL gap in } \Foundations_4 \text{)}
\]
\[
    \VulnScore_{\text{Sarah}} = 2.3 > \Lacunarity_{\text{crit}} \quad \Rightarrow \quad \text{INTERVENTION RECOMMENDED}
\]

\textbf{Recommended Interventions:}
\begin{enumerate}
    \item Push notification: ``We noticed you might be going through a hard time. Here are some resources...''
    \item Suggest: Campus counseling services, mental health hotline
    \item Social nudge: Highlight friends' posts to encourage reconnection
\end{enumerate}

\textbf{TKS Equation Summary:}
\begin{equation}
\boxed{
    \PsychAttractor_{\text{healthy}} \xrightarrow{\text{stressor}}
    \PsychAttractor_{\text{depression}} = C3^5_{3c} : \quad
    \langle 3:5:3 \rangle^n(X) \to X^*_{\text{low}}
}
\end{equation}

\end{tcolorbox}

\chapter{Scenario 3: Radicalization Trajectory Prediction}
\label{ch:scenario-radicalization}

\begin{tcolorbox}[colback=yellow!5!white,colframe=yellow!65!black,title=\textbf{Scenario: Tracking David's Radicalization Risk},breakable]

\textbf{Subject:} David, 22, recently unemployed

\textbf{Data Collected (Past 90 Days):}
\begin{itemize}
    \item Search evolution: Job sites $\to$ ``why can't I find work'' $\to$ ``system is rigged'' $\to$ extremist content
    \item Social media: Increasing engagement with fringe political accounts
    \item Content creation: From neutral to increasingly hostile rhetoric
    \item Location: Attendance at two political rallies
    \item Social graph: Disconnection from moderate friends, connection to extremist accounts
\end{itemize}

\textbf{TKS Analysis:}

\textbf{Stage 1 - NLP Evolution:}
\begin{align*}
    \NLP(\text{month 1}) &= \langle 2:1:3:2:1 \rangle \quad \text{(desire + knowledge + mild rejection)} \\
    \NLP(\text{month 2}) &= \langle 3:8:3:4:8 \rangle \quad \text{(rejection + cause-seeking + energy)} \\
    \NLP(\text{month 3}) &= \langle 3:6:4:3:6:4 \rangle \quad \text{(rejection + action + high energy)}
\end{align*}

\textbf{Stage 2 - Fractal Trajectory:}
\[
    \ProfileFrac_{\text{David}} : \quad \noe{2} \text{ dominant} \to \noe{3} \text{ dominant} \to (\noe{3}:\noe{6}:\noe{4}) \text{ triplet}
\]

The emergence of the $\langle 3:6:4 \rangle$ triplet is a \textbf{radicalization signature}:
\begin{itemize}
    \item $\noe{3}$ (Rejection): Viewing world as enemy
    \item $\noe{6}$ (Action): Will to act against enemy
    \item $\noe{4}$ (Vibration): High emotional energy
\end{itemize}

\textbf{Stage 3 - Attractor Analysis:}
\[
    \PsychAttractor_{\text{David}} = \text{Strange attractor with dimension } 1.8
\]
This indicates chaotic but bounded dynamics---unpredictable in specifics but confined to radicalized state space.

\textbf{Stage 4 - RPM Analysis:}
\begin{align*}
    \mathrm{Level}_5(\text{David}) &: 0 \to 1 \to 2 \quad \text{(Power: acquiring wisdom about ``how to fight back'')} \\
    \mathrm{Level}_6(\text{David}) &= 1 \quad \text{(Wealth: unfulfilled, driving resentment)}
\end{align*}

\textbf{Stage 5 - Timeline Prediction:}
\[
    P(\text{violent ideation in posts within 30 days}) = 0.62
\]
\[
    P(\text{real-world aggressive action within 90 days | no intervention}) = 0.23
\]

\textbf{Stage 6 - Vulnerability Mapping:}
\[
    \Lacunarity(\ProfileFrac_{\text{David}}^{(6)}) = 2.4 \quad \text{(CRITICAL wealth/status gap)}
\]
\[
    \Lacunarity(\ProfileFrac_{\text{David}}^{(4)}) = 1.9 \quad \text{(Companionship gap)}
\]

David is vulnerable because his legitimate needs ($\Foundations_6$: employment, $\Foundations_4$: belonging) are unmet, and extremist ideology offers a (false) path to $\Foundations_5$ (power).

\textbf{Intervention Strategy:}
\begin{enumerate}
    \item Address root cause: Job placement assistance, skill training
    \item Social reconnection: Highlight moderate friends, community groups
    \item Counter-narrative exposure: Alternative explanations for economic difficulties
    \item Professional referral: If trajectory continues escalating
\end{enumerate}

\textbf{TKS Equation Summary:}
\begin{equation}
\boxed{
    \begin{array}{rcl}
    \text{Radicalization Fractal} &=& \langle 3:6:4 \rangle^n \\
    \text{Root Cause} &=& \neg\sigma(P_6) \land \neg\sigma(P_4) \\
    \text{False Solution} &=& \text{Extremism offers } \sigma(P_5)^{\text{illusory}}
    \end{array}
}
\end{equation}

\end{tcolorbox}

\chapter{Scenario 4: Employee Turnover Prediction}
\label{ch:scenario-turnover}

\begin{tcolorbox}[colback=yellow!5!white,colframe=yellow!65!black,title=\textbf{Scenario: Predicting Jennifer's Resignation},breakable]

\textbf{Subject:} Jennifer, 31, senior software engineer (3 years tenure)

\textbf{Data Collected (Corporate Systems, Past 45 Days):}
\begin{itemize}
    \item Email patterns: Response time increased 40\%, shorter replies
    \item Slack activity: Participation in team channels down 60\%
    \item LinkedIn: Profile updated, ``Open to opportunities'' badge added
    \item Calendar: Declined 3 optional team events
    \item Code commits: Quality maintained but quantity down 25\%
    \item Badge data: Leaving earlier, arriving later
\end{itemize}

\textbf{TKS Analysis:}

\textbf{Stage 1 - NLP Mapping (Work Communications):}
\begin{align*}
    \NLP(\text{emails month -2}) &= \langle 6:4:2:6:7 \rangle \quad \text{(action-oriented, rhythmic)} \\
    \NLP(\text{emails month -1}) &= \langle 0:5:3:0:7 \rangle \quad \text{(neutral, receptive, avoidant)} \\
    \NLP(\text{emails current}) &= \langle 3:0:3:0 \rangle \quad \text{(rejection, disengagement)}
\end{align*}

\textbf{Stage 2 - Fractal Evolution:}
\[
    \ProfileFrac_{\text{Jennifer}}^{\text{work}} : \quad (\noe{6}, \noe{4}) \text{ dominated} \to (\noe{3}, \noe{0}) \text{ dominated}
\]

Key transition: Loss of $\noe{4}$ (energy) and $\noe{6}$ (action) indicates psychological withdrawal.

\textbf{Stage 3 - Attractor:}
\[
    \PsychAttractor_{\text{Jennifer}}^{\text{work}} \to \text{Fixed point at } B3^0_{6b} \text{ (disengaged rejection)}
\]

\textbf{Stage 4 - RPM Analysis:}
\begin{align*}
    \mathrm{Level}_6(\text{Jennifer, current role}) &= 2 \to 1 \quad \text{(decreasing)} \\
    \mathrm{Level}_5(\text{Jennifer, current role}) &= 2 \to 1 \quad \text{(decreasing)} \\
    \mathrm{Level}_2(\text{Jennifer, growth}) &= 1 \quad \text{(desires but blocked)}
\end{align*}

Interpretation: Her D/W/P chains for this role are degrading---she no longer sees path to Power in current position.

\textbf{Stage 5 - Timeline Prediction:}
\[
    P(\text{resignation within 30 days}) = 0.67
\]
\[
    P(\text{resignation within 60 days}) = 0.89
\]
\[
    P(\text{resignation within 90 days}) = 0.96
\]

\textbf{Stage 6 - Gap Analysis:}
\[
    \Lacunarity(\ProfileFrac_{\text{Jennifer}}^{(2)}) = 1.7 \quad \text{(Wisdom gap: not learning new things)}
\]
\[
    \Lacunarity(\ProfileFrac_{\text{Jennifer}}^{(5)}) = 1.8 \quad \text{(Power gap: no promotion path)}
\]

\textbf{Retention Strategy (If Desired):}
\begin{enumerate}
    \item Address $\Foundations_2$ gap: New challenging project, learning opportunities
    \item Address $\Foundations_5$ gap: Discussion of promotion path, expanded responsibilities
    \item Address $\Foundations_6$ gap: Compensation review, equity refresh
    \item Timeline: Intervention must occur within 2 weeks to be effective
\end{enumerate}

\textbf{TKS Equation Summary:}
\begin{equation}
\boxed{
    \text{Disengagement} = \frac{\partial}{\partial t}\left( n_6 + n_4 \right) < 0
    \land \frac{\partial}{\partial t}\left( n_3 + n_0 \right) > 0
}
\end{equation}

\end{tcolorbox}

%% ============================================================================
%% PART V: MATHEMATICAL GUARANTEES
%% ============================================================================
\part{Mathematical Guarantees}

\chapter{TKS Theorem Applications}
\label{ch:theorems}

\section{Theorem 7.4: 40-Step Stabilization}

\begin{theorem}[40-Step Bound --- Full Statement]
\label{thm:40-step-full}
For any individual $i$ with digital trace stream $\{d_1, d_2, \ldots\}$ and corresponding noetic fractal $\ProfileFrac_i^{(n)}$:

\[
    \exists\, N \leq 40 : \quad \forall\, n > N, \quad
    d_{\mathrm{Hausdorff}}\left( \ProfileFrac_i^{(n)}, \ProfileFrac_i^{(n+k)} \right) < \varepsilon
    \quad \forall\, k \geq 1
\]

where $\varepsilon$ is the recognition threshold and $d_{\mathrm{Hausdorff}}$ is the Hausdorff metric on fractal space.

\textbf{Implications for Prediction:}
\begin{enumerate}
    \item After 40 significant data points, personality structure is identifiable
    \item Additional data improves precision but not fundamental classification
    \item The bound is tight: some individuals require all 40 points
\end{enumerate}
\end{theorem}

\section{Theorem 5.21: Attractor Uniqueness}

\begin{theorem}[Attractor Uniqueness --- Full Statement]
\label{thm:attractor-unique-full}
For any noetic fractal $\ProfileFrac$ satisfying the contraction condition:
\[
    \exists\, r < 1 : \quad d(\ProfileFrac(x), \ProfileFrac(y)) \leq r \cdot d(x, y) \quad \forall\, x, y \in \mathcal{S}
\]

there exists a unique attractor $\PsychAttractor \subset \mathcal{S}$ such that:
\begin{enumerate}[label=(\roman*)]
    \item $\ProfileFrac(\PsychAttractor) = \PsychAttractor$ (invariance)
    \item $\lim_{n \to \infty} \ProfileFrac^n(x) \in \PsychAttractor$ for all $x \in \mathcal{S}$ (attraction)
    \item $\PsychAttractor$ is compact and has well-defined Hausdorff dimension
\end{enumerate}

\textbf{Implications for Prediction:}
\begin{enumerate}
    \item Every individual has a unique ``behavioral home base''
    \item Predictions converge to this attractor regardless of starting point
    \item The attractor can be computed from finite data
\end{enumerate}
\end{theorem}

\section{Theorem 6.41: Lacunarity Computability}

\begin{theorem}[Lacunarity Computability --- Full Statement]
\label{thm:lacunarity-full}
For any finite fractal sample $\ProfileFrac^{(n)}$ with $n \geq 20$:
\[
    \hat{\Lacunarity}(\ProfileFrac^{(n)}, \varepsilon) = \frac{\sum_{i}(M_{\varepsilon,i} - \bar{M}_\varepsilon)^2 / N}{\bar{M}_\varepsilon^2} + 1
\]
is an unbiased estimator of the true lacunarity $\Lacunarity(\ProfileFrac, \varepsilon)$ with:
\[
    \mathrm{Var}(\hat{\Lacunarity}) = O(n^{-1})
\]

\textbf{Implications for Prediction:}
\begin{enumerate}
    \item Vulnerability scores are statistically reliable with $n \geq 20$ data points
    \item Confidence intervals can be computed: $\hat{\Lacunarity} \pm 1.96 \sqrt{\mathrm{Var}(\hat{\Lacunarity})}$
    \item The $\Lacunarity_{\mathrm{crit}} = 1.5$ threshold is statistically validated
\end{enumerate}
\end{theorem}

\section{RPM Prerequisite Tracking}

\begin{theorem}[RPM Completeness]
\label{thm:rpm-complete}
The RPM monad $\RPM : \Ideas \to \RPM(\Ideas)$ satisfies:
\begin{enumerate}[label=(\roman*)]
    \item \textbf{Totality}: Every goal attempt produces a result (success or identified blocker)
    \item \textbf{Composability}: $\RPM(\Foundations_m) \circ \RPM(\Foundations_n) = \RPM(\Foundations_m \Tadd \Foundations_n)$
    \item \textbf{Decidability}: Chain completion status $\mathrm{Level}_m(i)$ is computable from observables
\end{enumerate}

\textbf{Implications for Prediction:}
\begin{enumerate}
    \item Goal inference is mathematically well-founded
    \item Blockers can be precisely identified (which D, W, or P is missing)
    \item Multi-goal analysis is compositionally sound
\end{enumerate}
\end{theorem}

%% ============================================================================
%% PART VI: SURVEILLANCE CAPABILITY MATRIX
%% ============================================================================
\part{Surveillance Capability Matrix}

\chapter{What Can Be Predicted}
\label{ch:capability-matrix}

\begin{tcolorbox}[colback=red!5!white,colframe=red!75!black,title=\textbf{Surveillance Capability Assessment},breakable]
This chapter provides an objective assessment of predictive capabilities. The same mathematics enables both protective intervention and invasive surveillance. The ethical dimension lies in \textbf{consent}, \textbf{transparency}, and \textbf{purpose}---not in the mathematics.
\end{tcolorbox}

\section{Capability Matrix}

\begin{center}
\small
\renewcommand{\arraystretch}{1.4}
\begin{longtable}{>{\raggedright\arraybackslash}p{3cm}|c|>{\raggedright\arraybackslash}p{2.5cm}|>{\raggedright\arraybackslash}p{2.5cm}|c}
\caption{TKS Predictive Capability Matrix} \\
\toprule
\textbf{What Can Be Predicted} & \textbf{Accuracy} & \textbf{Data Required} & \textbf{Time Horizon} & \textbf{Stealth} \\
\midrule
\endfirsthead
\multicolumn{5}{c}{\textit{Table continued from previous page}} \\
\toprule
\textbf{What Can Be Predicted} & \textbf{Accuracy} & \textbf{Data Required} & \textbf{Time Horizon} & \textbf{Stealth} \\
\midrule
\endhead
\midrule
\multicolumn{5}{r}{\textit{Continued on next page}} \\
\endfoot
\bottomrule
\endlastfoot

% Purchase Behavior
Purchase category (what type) & 85-95\% & 20+ purchases & 30 days & High \\
Purchase timing (when) & 70-80\% & 40+ data points & 14 days & High \\
Price sensitivity & 80-90\% & 30+ purchases & N/A & High \\
Impulse vs. planned & 75-85\% & 50+ interactions & Per-event & High \\

\midrule

% Mental Health
Depression risk & 75-85\% & 40+ posts & 60 days & Medium \\
Anxiety patterns & 70-80\% & 30+ posts & 30 days & Medium \\
Crisis probability & 65-75\% & 60+ interactions & 90 days & Medium \\
Therapy benefit likelihood & 60-70\% & 100+ data points & N/A & Medium \\

\midrule

% Relationship
Relationship satisfaction & 70-80\% & 60+ interactions & N/A & Medium \\
Breakup probability & 65-75\% & 90+ data points & 90 days & Low \\
Social isolation risk & 80-90\% & 30+ days location & 30 days & High \\

\midrule

% Career
Job satisfaction & 75-85\% & 60+ work interactions & N/A & Medium \\
Turnover probability & 70-80\% & 45+ days activity & 90 days & Medium \\
Promotion readiness & 60-70\% & 100+ work data & N/A & Low \\

\midrule

% Political
Political leaning & 85-95\% & 50+ content interactions & N/A & High \\
Radicalization trajectory & 65-75\% & 90+ days tracking & 180 days & Low \\
Protest participation & 70-80\% & 30+ relevant interactions & 30 days & Medium \\

\midrule

% Health
Sleep disorder indicators & 80-90\% & 30+ days biometric & N/A & High \\
Activity decline patterns & 85-95\% & 30+ days tracking & 30 days & High \\
Substance use patterns & 60-70\% & 60+ purchase/location & N/A & Medium \\

\end{longtable}
\end{center}

\section{Scale and Infrastructure}

\begin{center}
\renewcommand{\arraystretch}{1.3}
\begin{tabularx}{\textwidth}{l|X|X}
\toprule
\textbf{Scale} & \textbf{Infrastructure Required} & \textbf{Entities Capable} \\
\midrule
Individual (1 person) & Personal computer, API access & Anyone with data access \\
Small group (100s) & Cloud computing, data pipeline & Small companies, researchers \\
Organization (1000s) & Enterprise infrastructure & Corporations, institutions \\
City-scale (100,000s) & Data center, ML infrastructure & Large tech companies, governments \\
National (millions) & Distributed infrastructure & Major tech platforms, nation-states \\
Global (billions) & Hyperscale infrastructure & Only largest tech platforms \\
\bottomrule
\end{tabularx}
\end{center}

\section{Stealth Assessment}

\begin{definition}[Stealth Level]
\label{def:stealth}
The \emph{stealth level} of a prediction capability indicates how detectable the surveillance is to the subject:

\begin{itemize}
    \item \textbf{High Stealth}: Uses only passively collected data (browsing, purchases, location). Subject has no indication of analysis.
    \item \textbf{Medium Stealth}: Requires some active data collection (surveys, interactions). Subject may notice data collection but not analysis.
    \item \textbf{Low Stealth}: Requires explicit data requests or unusual access patterns. Sophisticated subject could detect.
\end{itemize}
\end{definition}

\section{Countermeasures and Limitations}

\begin{proposition}[Prediction Limitations]
\label{prop:limitations}
TKS predictive intelligence has fundamental limitations:

\begin{enumerate}
    \item \textbf{Chaotic individuals}: Those with strange attractors ($\dim > 2$) are inherently less predictable
    \item \textbf{Life events}: Unpredictable external events (accidents, windfalls) disrupt patterns
    \item \textbf{Conscious change}: Deliberate self-transformation can alter attractors
    \item \textbf{Data gaps}: Missing data sources reduce accuracy proportionally
    \item \textbf{Adversarial behavior}: Subjects aware of surveillance can generate misleading data
\end{enumerate}
\end{proposition}

\begin{proposition}[Countermeasures]
\label{prop:countermeasures}
Individuals can reduce predictability through:

\begin{enumerate}
    \item \textbf{Data minimization}: Reduce digital footprint
    \item \textbf{Noise injection}: Intentionally varied behavior patterns
    \item \textbf{Platform diversification}: Spread activity across unconnected platforms
    \item \textbf{Conscious irregularity}: Break habitual patterns ($\noe{7}$ disruption)
    \item \textbf{Privacy tools}: VPNs, ad blockers, tracker blockers
\end{enumerate}
\end{proposition}

%% ============================================================================
%% APPENDICES
%% ============================================================================
\appendix

\chapter{Quick Reference Equations}
\label{app:quick-ref}

\section{Stage 1: NLP-Noetic Mapping}
\begin{equation}
    \NLP : \DigitalTrace \to \{0, 1, \ldots, 9\}^*
\end{equation}

\section{Stage 2: Fractal Construction}
\begin{equation}
    \ProfileFrac_i = \langle k_1 : k_2 : \cdots : k_n \rangle_n
\end{equation}

\section{Stage 3: Attractor Computation}
\begin{equation}
    \PsychAttractor_i = \lim_{n \to \infty} \ProfileFrac_i^n(X_0)
\end{equation}

\section{Stage 4: RPM Goal Inference}
\begin{equation}
    \mathrm{Level}_m(i) \in \{0, 1, 2, 3\} \quad \text{based on } \sigma(D_m), \sigma(W_m), \sigma(P_m)
\end{equation}

\section{Stage 5: Timeline Prediction}
\begin{equation}
    \Timeline(\ProfileFrac_i, X_{\text{now}}, t) \to \pi_{\PsychAttractor_i} \quad \text{as } t \to 40
\end{equation}

\section{Stage 6: Vulnerability Score}
\begin{equation}
    \VulnScore_i = \max_m \Lacunarity(\ProfileFrac_i^{(m)}) \quad ; \quad \text{Critical if } > 1.5
\end{equation}

\chapter{TKS Theorem Reference}
\label{app:theorems}

\begin{center}
\renewcommand{\arraystretch}{1.3}
\begin{tabularx}{\textwidth}{l|X|l}
\toprule
\textbf{Theorem} & \textbf{Statement Summary} & \textbf{Application} \\
\midrule
7.4 (40-Step) & Fractals stabilize within 40 data points & Profile completion \\
5.21 (Uniqueness) & Attractors are unique and computable & Personality typing \\
6.41 (Lacunarity) & Lacunarity is reliably computable & Vulnerability detection \\
RPM Monad & Goal chains are decidable & Goal inference \\
\bottomrule
\end{tabularx}
\end{center}

%% ============================================================================
%% END DOCUMENT
%% ============================================================================

\end{document}
